% !TeX program = xelatex
% SPDX-License-Identifier: GPL-3.0-only

% To use a project copy of SimSun, uncomment the next line and match its path and filename.
% \AddToHook{fduthesis/cjk-font}{\setCJKmainfont[Path=fonts/,FontIndex=0,AutoFakeBold=2]{simsun.ttc}}
\documentclass[lang=en]{fduthesis}

\usepackage{booktabs}

\usepackage[backend=biber,style=gb7714-2025]{biblatex}
\addbibresource{reference.bib}

\fdusetup{
  % Uncomment the next line to number the main text i, ii, iii instead of 1, 2, 3.
  % main-page-numbering = roman,
  %
  degree = master, type = academic, clc = TP3,
  %
  title = {顺序查找何时优于二分查找？},
  title-en = {When Does a Linear Scan Beat Binary Search?},
  %
  author = {张三}, student-id = {123456789},
  % Use 同等学力人员 for equivalent qualifications, or none otherwise.
  % With lang=en, the cover adds 英文项目 automatically.
  student-category = none,
  %
  supervisor = {李四}, supervisor-title = {教授},
  committee = {李四教授, 王五副教授},
  %
  major = {计算机科学与技术}, department = {计算与智能创新学院},
  date = {2026年6月30日},
  %
  keywords = {二分查找, 有序数组, 查询频率},
  keywords-en = {binary search, sorted arrays, query frequency}
}

\begin{document}
\frontmatter
\makecoverpage
\makecommittee
\tableofcontents
\begin{fdusummary}
在按姓名排序的通讯录中查找一个人，可以从第一项开始逐项核对，也可以先看中间的一项，再根据比较结果只查前半段或后半段。二分查找能限制一次查询最多需要访问多少项，但如果经常要找的人恰好排在第一位，顺序查找往往只需查看一次。本文研究一个可以逐项核算的问题：在同一组有序数据上，第一项需要多常被查到，顺序查找的平均访问次数才会低于二分查找？为单独考察查询频率的影响，分析过程中固定数据的内容、排列次序和两种算法的访问规则，仅改变目标落在不同位置的概率。

研究对象是包含七个互异键的升序数组，首先考虑每次查询的目标都在数组中的情形。访问一个位置并判断目标与该位置的键之间的大小关系，计为一次探查。顺序查找总是从首项开始，命中便停止。二分查找从第四项开始，根据比较结果选择左侧或右侧的三个候选位置，随后继续取中点。以查找第五项为例，算法先查看第四项，接着查看第六项，最后才在第五项命中。虽然第四项和第五项在数组中相邻，查到它们所需的探查次数分别为一和三。由此可见，数组位置对于两种算法有不同的含义，不能只按键与数组中点的距离估算二分查找的代价。

本文先按查询目标列出两种算法的探查次数，再以查询概率对各列加权。七个位置等概率被查询时，顺序查找的探查次数从一至七逐项增加，平均为四次。二分查找有一个位置可在第一次命中，两个位置可在第二次命中，其余四个位置在第三次命中，平均为七分之十七次。这个基准同时给出了两种算法各自的最坏情况。只看最坏情况会选择二分查找，但这还不能回答在同一张通讯录上反复查询时应当采用哪种方法。

随后令第一项被查询的概率为 $p$，把剩余概率等分给另外六项。查询第一项时，顺序查找只探查一次，二分查找需要三次。对其余六项各查一次，顺序查找共需二十七次，二分查找共需十四次。因此，提高首项查询概率有两种同时发生的效果：顺序查找获得更多一次命中的机会，二分查找则要更频繁地走到判定树的第三层。两种平均代价随首项概率呈线性变化，变化方向相反。解得交点为 $p=13/25$，此时两种算法平均均探查 $67/25$ 次。低于这一概率时二分查找占优，高于这一概率时顺序查找占优。

为核对加权过程，本文将等概率查询和只查第一项分别代入公式。前者得到与逐键求和相同的平均值，后者得到顺序查找一次、二分查找三次，与实际查找路径一致。在交点处，两种算法的平均值相同，却没有一个位置需要恰好探查这么多次。这个分数表示按给定频率进行多次查询后，每次查询分摊到的探查次数。例如，在二十五次查询中安排十三次查询首项，其余六项各查两次，顺序查找和二分查找都刚好探查六十七次。把其中一次首项查询移到第四项，顺序查找便增加三次探查，二分查找却减少两次。这个例子说明，预测哪种算法更合适还需要知道热点落在哪个位置。

将数组扩展到十五个键，并保持首项以外的成功查询等概率时，交点升至 $37/58$。回到七个键，若允许查询不存在的键，顺序查找遇到第一个更大的键即可停止。假设八个空档等概率，在成功查询中又有四分之三查询首项，则失败查询占全部请求的比例达到 $23/56$ 时，两种算法再次持平。将方法用于其他有序表时，可以先记录两种查找方式访问每个目标所需的次数，再按实际查询分布计算期望。这样可以看出哪些位置的查询左右了算法选择，也明确了由探查次数比较结果推断运行时间时，还需要另外测量每次访问的实际开销。
\end{fdusummary}

\begin{fdusummaryen}
Binary search limits the number of probes needed to find any key in a sorted
array, but a linear scan can find a frequently requested first key in a single
probe. We compare successful searches on the same seven distinct keys, counting
one inspection of an array element as one probe. With uniform queries, a linear
scan averages four probes and binary search averages $17/7$. We then assign
probability $p$ to the first key and divide the remaining probability equally
among the other six. A linear scan takes one probe for the first key and
a total of 27 probes across one search for each of the other six keys. The corresponding counts for
binary search are three and 14. The expected costs coincide at $p=13/25$,
when both methods average $67/25$ probes. For example, both need exactly
67 probes for 13 requests for the first key and two for each remaining key. Above that threshold the linear scan
uses fewer probes, although binary search still has the smaller worst-case
cost. Substituting $p=1/7$ and $p=1$ recovers the uniform-query averages and
the costs of finding the first key. The threshold depends on the location of
the frequent key and the distribution of the remaining queries. Increasing the
array to fifteen keys raises the first-key threshold to $37/58$. In the original
array, suppose three quarters of successful queries ask for the first key and
misses fall uniformly among the eight gaps. The linear scan then wins while
fewer than $23/56$ of all queries miss.
\end{fdusummaryen}
\listoffigures
\listoftables

\mainmatter
\chapter{Counting probes}
Suppose a directory stores seven names in alphabetical order. A linear scan
starts at the top of the list. Binary search starts at the fourth name, then
searches the half that could still contain the target. We count one probe each
time a search inspects a name, following the comparison model used for ordered
search by Knuth\cite{knuth1998}.

\begin{table}[htbp]
  \centering
  \caption{Probes needed to find each key}\label{tab:counts}
  \begin{tabular}{lrrrrrrr}
    \toprule
    Key position & 1 & 2 & 3 & 4 & 5 & 6 & 7 \\
    \midrule
    Linear scan   & 1 & 2 & 3 & 4 & 5 & 6 & 7 \\
    Binary search & 3 & 2 & 3 & 1 & 3 & 2 & 3 \\
    \bottomrule
  \end{tabular}
\end{table}

The path in Figure~\ref{fig:path} explains the entry for key 5 in
Table~\ref{tab:counts}. The first probe rules out keys 1 through 4.
The second leaves a single candidate, which the third probe confirms.

\begin{figure}[htbp]
  \centering
  \fbox{Probe 4}\quad$\longrightarrow$\quad
\fbox{Probe 6}\quad$\longrightarrow$\quad\fbox{Find 5}
%
  \caption{Binary search for the fifth key}\label{fig:path}
\end{figure}

\section{What one probe measures}
A search that visits positions 4, 6 and 5 costs three probes even if the
implementation tests equality and order in separate instructions. The model
counts the records that must be inspected. It is useful, for example, when a
record has to be fetched before either test can proceed. Translating a probe
count into elapsed time would also require the cost of fetching that record.
A scan reads consecutive records, whereas the midpoint rule can jump between
distant locations.

For binary search, the candidates always occupy a closed interval. If the
middle key is too small, every key to its left is also too small, so the next
interval begins immediately to its right. At the end of the path in
Figure~\ref{fig:path}, the interval consists of position 5 alone. It must
still be inspected. Stopping as soon as the two endpoints coincide would
turn this successful search into a miss.

\chapter{Changing the query distribution}
Let $p$ be the probability that a query asks for the first key. Each of the
other six keys has probability $(1-p)/6$. Summing the last six columns of
Table~\ref{tab:counts} gives 27 probes for the linear scan and 14 for binary
search. Their expected costs are therefore
\begin{equation}
  L(p)=p+27(1-p)/6, \qquad B(p)=3p+14(1-p)/6.
  \label{eq:cost}
\end{equation}
Solving $L(p)=B(p)$ gives $p=13/25$. Above that threshold the linear scan
wins. At the threshold both searches average $67/25$ probes. A worst-case
bound alone would miss this change in their ranking.

\section{Accounting for a batch of requests}
The break-even distribution has a convenient whole-number interpretation.
Imagine 25 requests written on index cards: 13 for the first key and two
for each of the other six. The linear scan spends 13 probes on the first
pile and $2\times27=54$ on the rest. Binary search spends 39 probes on the
first pile and 28 on the rest. Both totals are 67, independently of the order
in which the cards are drawn.

Now change one card from key 1 to key 4. The linear scan spends three extra
probes on that request, while binary search saves two. The new totals are 70
and 65. Changing the card to key 2 instead produces totals of 68 and 66.
Both edits reduce the number of requests for key 1 by exactly one. Their
different effects show why the threshold from Equation~\ref{eq:cost} depends
on how requests for the remaining keys are distributed.

For an observed batch, let $c_i$ be the number of requests for key $i$.
Subtracting the two rows of Table~\ref{tab:counts} gives the total excess
cost of the scan over binary search:
\begin{equation}
  D=-2c_1+3c_4+2c_5+4c_6+4c_7.
  \label{eq:batch}
\end{equation}
The two methods tie on keys 2 and 3, so those requests contribute zero to
the difference. A negative $D$ favors the scan. This calculation can be
applied directly to counts from a directory log without first fitting a
probability distribution.

\chapter{A larger directory}
\section{Summing probes by depth}
Consider directories of size $n=2^h-1$. Each midpoint splits the remaining
keys into two groups of equal size. There is one key found on probe 1, two
more on probe 2, and $2^{j-1}$ keys first found on probe $j$. Write $S_h$
for the total binary-search cost of looking up every key once. Growing the
tree by a level adds $2^h$ keys at depth $h+1$, so
\begin{equation}
  S_1=1, \qquad S_{h+1}=S_h+(h+1)2^h .
  \label{eq:level}
\end{equation}
For seven keys this gives $S_3=17$. The next eight keys cost four probes
each, bringing the fifteen-key total to 49. The linear scan, by comparison,
uses the sum of the positions, $T_n=n(n+1)/2$, or 120 probes for fifteen keys.

\section{How much weight must the first key carry?}
On the first key the scan saves $h-1$ probes. On the remaining keys, looking
up each one once costs the scan $T_n-1$ probes and binary search $S_h-h$.
Let $R=(T_n-1)-(S_h-h)$ denote this excess. If the first key receives
probability $p$ and the other keys share the rest equally, the scan breaks
even when its savings on the first key pay for the excess on the rest:
\begin{equation}
  p(h-1)=(1-p)\frac{R}{n-1}.
  \label{eq:size-threshold}
\end{equation}

\begin{table}[htbp]
  \centering
  \caption{First-key frequency required for equal expected cost}
  \label{tab:sizes}
  \begin{tabular}{rrrrr}
    \toprule
    Keys $n$ & Levels $h$ & Scan total $T_n$ & Binary total $S_h$ & Threshold $p$ \\
    \midrule
    3  & 2 &   6 &   5 & $1/2$ \\
    7  & 3 &  28 &  17 & $13/25$ \\
    15 & 4 & 120 &  49 & $37/58$ \\
    31 & 5 & 496 & 129 & $371/491$ \\
    \bottomrule
  \end{tabular}
\end{table}

Suppose the first key gets three requests out of five. In the seven-key
directory that share is above $13/25$, but in the fifteen-key directory it
falls below $37/58$. The same access pattern therefore favors different
methods at these two sizes. Table~\ref{tab:sizes} uses complete trees so
that each cost can be checked from Equation~\ref{eq:level}. For an arbitrary
size, the same accounting works after recording the actual midpoint depths.

\chapter{When a requested key is absent}
\section{Counting the gaps}
Return to seven keys. An unsuccessful request lies in one of eight gaps:
before the first key, between consecutive keys, or after the last key.
A scan may stop at the first key larger than the target because all later
keys are larger still. A request between keys 2 and 3 therefore costs three
probes. Binary search for that request probes positions 4, 2 and 3, leaving
an empty interval. In this complete tree every unsuccessful path takes three
probes.

\begin{table}[htbp]
  \centering
  \caption{Cost of an unsuccessful search, by insertion position}
  \label{tab:gaps}
  \begin{tabular}{lrrrrrrrr}
    \toprule
    Keys smaller than target & 0 & 1 & 2 & 3 & 4 & 5 & 6 & 7 \\
    \midrule
    Linear scan   & 1 & 2 & 3 & 4 & 5 & 6 & 7 & 7 \\
    Binary search & 3 & 3 & 3 & 3 & 3 & 3 & 3 & 3 \\
    \bottomrule
  \end{tabular}
\end{table}

The last two columns of Table~\ref{tab:gaps} cost the scan the same amount.
Both require inspecting the final record. Passing the end of the directory
adds no record inspection. Under equally likely gaps the scan averages
$35/8$ probes per miss, an excess of $11/8$ over binary search.

\section{Combining hits and misses}
Let $q$ be the fraction of all requests that miss, and suppose three quarters
of the successful requests ask for the first key. The other six keys divide
the remaining successful requests equally. Substituting $p=3/4$ into
Equation~\ref{eq:cost} gives $15/8$ probes for a successful scan and $17/6$
for a successful binary search. The scan saves $23/24$ probes per hit. If
the eight gaps are equally likely conditional on a miss, the expected excess
cost of the scan per request is
\begin{equation}
  -(1-q)\frac{23}{24}+q\frac{11}{8}=\frac{56q-23}{24}.
  \label{eq:miss}
\end{equation}
At $q=1/4$ the scan saves $3/8$ probes per request. At $q=1/2$ it costs
an extra $5/24$. The crossing occurs at $q=23/56$. The location of misses
matters as much as their frequency: if all misses fall before the first key,
each miss saves the scan two probes, and both hits and misses now favor it.

\printbibliography

\appendix
\chapter{Checking the endpoints}
At $p=1/7$ all seven keys are equally likely. Equation~\ref{eq:cost}
gives $L=4$ and $B=17/7$, the row averages in Table~\ref{tab:counts}.
At $p=1$ every query asks for the first key, and the same equation gives
one and three probes respectively.

\section{An integer check on the mixed workload}
Equation~\ref{eq:miss} predicts a tie for 23 misses in every 56 requests.
To reconstruct a batch with whole counts in every gap and at every key, use
448 requests. Of these, 184 miss, giving 23 per gap. The other 264 succeed:
198 ask for the first key and 11 ask for each of the other six.
For the linear scan the hits cost $198+11\times27=495$ probes and the
misses cost $23\times35=805$. For binary search the hits cost
$198\times3+11\times14=748$ and the misses cost $184\times3=552$.
Both methods use exactly 1,300 probes. This check enumerates the requests
directly, so it also tests that the first-key probability was applied to
the successful requests alone.

\backmatter

\begin{acknowledgements}
I thank my supervisor for suggesting that I check both endpoints of the
probability model.
\end{acknowledgements}
\makeauthorization
\end{document}
